\documentclass[11pt,a4paper]{article}
\usepackage[margin=0.85in,top=1in,bottom=1in]{geometry}
\usepackage{amsmath,amssymb,amsfonts}
\usepackage{graphicx}
\usepackage{float}
\usepackage{enumitem}
\usepackage{microtype}
\usepackage{caption}
\usepackage[hidelinks]{hyperref}
\usepackage[most,breakable]{tcolorbox}
\tcbuselibrary{breakable,skins}
\usepackage{fontspec}
\setmainfont{DejaVu Serif}
\setsansfont{DejaVu Sans}
\setmonofont{DejaVu Sans Mono}
\captionsetup{font=small,labelfont=bf,skip=4pt}
\setlist{leftmargin=1.5em,topsep=2pt}

%% Breakable problem card — prevents content cutoff across pages
\newtcolorbox{solcard}[3]{
  breakable, enhanced,
  colback=white, colframe=gray!50, arc=2pt,
  title={\textbf{#1}\hfill\textsf{\small #3\,pts}},
  fonttitle=\normalsize\sffamily\bfseries,
  before upper={\small\textit{#2}\par\smallskip\hrule height 0.4pt\smallskip},
  top=4pt, bottom=6pt, left=7pt, right=7pt,
  parbox=false,
}

\title{\textbf{Transformers --- Solution}}
\author{\normalsize X26 (2026) — English translation}
\date{}
\begin{document}\maketitle\vspace{-0.5em}
\begin{solcard}{A1}{What are the magnetic field fluxes through the primary and secondary windings $\Phi_1$ and $\Phi_2$? Express your answers using $\Phi$, $N_1$ and $N_2$.}{0.20}
\par\smallskip\noindent\textbf{Answer:} \textbackslash{}textbf{Answer:} $$\Phi_1=\Phi \cdot N_1,\qquad \Phi_2 =- \Phi \cdot N_2$$
\end{solcard}

\begin{solcard}{A2}{Using the circulation theorem for the magnetic field strength vector $\vec H$, find the magnetic field flux $\Phi$ in the core. Express your answer in terms of $S$, $\ell$, $\mu$, $I_1$, $I_2$, $N_1$ and $N_2$.}{0.20}

Let us write the circulation theorem for the entire transformer: $$H\cdot l = N_1I_1-N_2I_2,$$Let us express the magnetic field flux:

\par\smallskip\noindent\textbf{Answer:} \textbackslash{}textbf{Answer:} $$\Phi=BS=\mu \mu_0HS= \frac{\mu\mu_0S}{l}(N_1I_1 - N_2I_2)$$
\end{solcard}

\begin{solcard}{A3}{Obtain expressions for the inductances of windings $L_1$, $L_2$ and their mutual inductance coefficient $M$. Express your answer using $S$, $\ell$, $\mu$, $I_1$, $I_2$, $N_1$ and $N_2$.}{0.30}
\par\smallskip\noindent\textbf{Answer:} \textbackslash{}textbf{Answer:} $$ \begin{aligned} & L_1=\frac{\partial \Phi_1}{\partial I_1}=\frac{\mu \mu_0 S N_1^2}{l} \\ & L_2=\frac{\partial \Phi_2}{\partial I_2}=\frac{\mu \mu_0 S N_2^2}{l} \\ & M=\frac{\partial \Phi_1}{\partial I_2}=\frac{\partial \Phi_2}{\partial I_1}=-\frac{\mu \mu_0 S N_1N_2}{l} \end{aligned} $$
\end{solcard}

\begin{solcard}{A4}{Find the magnetic induction emf $\mathcal E_{1}$ and $\mathcal E_{2}$ arising in the primary and secondary windings, respectively. Express your answer using $\dot{{I}_{1}}, \dot{{I}_{2}}, L_{1}, L_{2}$ and $M$. Notes: EMF $\mathcal E_{1}$ is considered positive if the electrostatic potential at the top point a1 is greater than at the bottom a2. EMF $\mathcal E_{2}$ is considered positive if the electrostatic potential at the top point b1 is greater than at the bottom b2.}{0.40}

Find the magnetic induction emf $\mathcal E_{1}$ and $\mathcal E_{2}$ arising in the primary and secondary windings, respectively. Express your answer using $\dot{{I}_{1}}, \dot{{I}_{2}}, L_{1}, L_{2}$ and $M$.


Notes:

\par\smallskip\noindent\textbf{Answer:} \textbackslash{}textbf{Answer:} $$ \begin{aligned} & \mathcal E_1=\dot\Phi_1=N_1 \dot{\Phi}=L_1 \dot{I}_1+M \dot{I}_2 \\ & \mathcal E_2=-\dot\Phi_2=N_2 \dot{\Phi}=-L_2 \dot{I}_2-M \dot{I}_1 \end{aligned} $$
\end{solcard}

\begin{solcard}{A5}{Obtain an expression for the current transformation ratio $k_{I}\equiv{\left|\tilde{I}_{1}\right|}/{\left|\tilde{I}_{2}\right|}$. The answer may include $L_1$, $L_2$, $\omega$ and $Z$.}{0.50}

Using the equations from A4 and the following relations: $$U_2 = I_2Z\\ M^2=L_1L_2$$ $$ I_2 Z=\mathcal E_2=-L_2 \dot{I}_2-M \dot{I}_1 \\ I_2\left(Z-i \omega L_2\right)=i \omega M I_1 $$ we obtain

\par\smallskip\noindent\textbf{Answer:} \textbackslash{}textbf{Answer:} $$ k_I=\left|\frac{\tilde{I}_1}{\tilde{I}_2}\right|=\left|\frac{Z-i \omega L_2}{i \omega M}\right|=\frac{\sqrt{(\operatorname{Re} Z)^2+\left(\text{Im}~ Z-\omega L_2\right)^2}}{\omega \sqrt{L_1 L_2}} $$
\end{solcard}

\begin{solcard}{A6}{What is the voltage transformation ratio $k_{U}={\left|\tilde{U}_{2}\right|}/{\left|\tilde{U}_{1}\right|}$? The answer may include $L_1$, $L_2$, $\omega$ and $Z$.}{0.50}

Writing down the equations from the previous paragraphs: $$ U_2=-L_2 \dot{I}_2-M \dot{I}_1 \\ U_2=I_2 \cdot Z \\ U_1=L_1 \dot{I}_1+M \dot{I}_2 $$ we express $U_2$ via/through $U_1$: $$ U_2=i \omega L_2 \frac{U_2}{Z}+i \omega M I_1 \\ U_1=-i \omega L_1 I_1-i \omega M \frac{U_2}{Z} \\ U_2\left(1-\frac{i \omega L_2}{Z}\right)=\frac{M}{L_1}\left(U_1-\frac{i \omega M}{Z} U_2\right)\\ \frac{U_2}{U_1}=\frac{\frac{M}{L_1}}{1-\frac{i \omega L_2}{Z}+\frac{i \omega M^2}{Z L_1}}=\frac{M}{L_1}=-\frac{N_2}{N_1} $$ we get:

\par\smallskip\noindent\textbf{Answer:} \textbackslash{}textbf{Answer:} $$ k_U=\left|-\frac{N_2}{N_1}\right|=\sqrt{\dfrac{L_2}{L_1}} $$
\end{solcard}

\begin{solcard}{A7}{Write down the condition for the number of turns in windings $N_{1}$ and $N_{2}$, under which the expression $(1)$ is satisfied.}{0.30}
\par\smallskip\noindent\textbf{Answer:} \textbackslash{}textbf{Answer:} $$ N_2 \gg N_1 $$
\end{solcard}

\begin{solcard}{A8}{Determine the phase difference $\Delta\varphi$ by which voltage $U_1$ leads current $I_1$. The answer may include $L_1$, $L_2$, $\omega$ and $Z$.}{0.60}

$$ \begin{aligned} & U_1+\frac{i \omega M}{Z} \cdot U_2=-i \omega L_1 I_1 \\ & U_1\left(1-\frac{i \omega M N_2}{Z N_1}\right)=-i \omega L_1 I_1\\ & U_1 = \frac{-i\omega L_1 Z N_1}{ZN_1 - i\omega MN_2}I_1 \end{aligned} $$after transformations:


Or, which is convenient for further calculations, $$\cos{\Delta \varphi} = \frac{\omega L_2 \text{Re}~Z}{|Z|\sqrt{(\text{Re}~Z)^2+(\omega L_2-\text{Im}~Z)^2}}$$

\par\smallskip\noindent\textbf{Answer:} \textbackslash{}textbf{Answer:} $$\text{tg}{\Delta\varphi}=\frac{|Z|^2- \omega L_2 \cdot \text{Im}~Z}{\omega L_2 \text{Re}~Z}$$
\end{solcard}

\begin{solcard}{А9}{Determine the time-averaged power $P$ of source $U_1$. The answer may include $U_0$, $Z$, $N_1$ and $N_2$.}{0.60}

The average power is found as: $$P=\frac{1}{2} |I_1| |U_1|\cos{\Delta \varphi}$$Substituting we get:


The calculations could be significantly reduced by using the following formula for power: $$P=\text{Re}\left(\frac{1}{2} I_1 U_1^*\right)$$

\par\smallskip\noindent\textbf{Answer:} \textbackslash{}textbf{Answer:} $$P = \frac{U_1^2 \text{Re}~Z}{2 |Z|^2}\cdot \frac{N_2^2}{N_1^2}$$
\end{solcard}

\begin{solcard}{B1}{Find the ratio $I / I_{A}$.}{0.50}

Using the formula for $k_I$ we get the answer:

\par\smallskip\noindent\textbf{Answer:} \textbackslash{}textbf{Answer:} $$\frac{I}{I_A}=\sqrt{\frac{L_1}{L_2}} = N$$
\end{solcard}

\begin{solcard}{C1}{Obtain an expression for the average power $P_0$, which is released at the load if the winding resistance can be neglected. Express your answer through the amplitude of the input voltage $U_0$ and the load resistance $R$. Calculate it for a transformer whose parameters were given in the problem earlier.}{0.20}

Using the formula for average power:

\par\smallskip\noindent\textbf{Answer:} \textbackslash{}textbf{Answer:} $$P_0=\frac{U_0^2}{2R} = \frac{U_\mathrm{rms}^2}{R} = 303 ~\text{W}$$
\end{solcard}

\begin{solcard}{C2}{Similar to Part A, write equations relating the complex amplitudes $\tilde U_1$, $\tilde U_2$, $\tilde I_1$ and $\tilde I_2$. The response may also include $R$, $r$, single winding inductance $L$, and cyclic frequency $\omega$.}{0.30}
\par\smallskip\noindent\textbf{Answer:} \textbackslash{}textbf{Answer:} $$\tilde{U_1}=-i \omega L \tilde{I}_1+i \omega L \tilde{I}_2+r \tilde{I}_1 \\ \tilde{U}_2=(R+r) \tilde{I}_2 \\ \tilde{U}_2=i \omega L \tilde{I}_2-i \omega L \tilde{I}_1 \\$$
\end{solcard}

\begin{solcard}{C3}{Express $\tilde I_1$ and $\tilde I_2$ in terms of $\tilde U_1$ and other quantities you need, which may include $R$, $r$, $L$ and $\omega$.}{0.40}

We express: \textbackslash{}begin{gathered} \textbackslash{}tilde{I}_2(R+r-i \textbackslash{}omega L)=-i \textbackslash{}omega L \textbackslash{}tilde{I}_1 \textbackslash{}\textbackslash{} \textbackslash{}tilde{U}_1=\textbackslash{}left(-i \textbackslash{}omega L+r+\textbackslash{}frac{\textbackslash{}omega^2 L^2}{R+r-i \textbackslash{}omega L}\textbackslash{}right) \textbackslash{}tilde{I}_1 \textbackslash{}\textbackslash{} \textbackslash{}tilde{I}_1=\textbackslash{}frac{(R+r-i \textbackslash{}omega L) \textbackslash{}tilde{U}_1}{(r-i \textbackslash{}omega L)(R+r-i \textbackslash{}omega L)+\textbackslash{}omega^2 L^2} \textbackslash{}end{gathered} we get:

\par\smallskip\noindent\textbf{Answer:} \textbackslash{}textbf{Answer:} $$\tilde{I}_1=\frac{R+r-i \omega L}{r(R+r)-i \omega L(R+2 r)} \tilde{U}_1 $$$$ \tilde{I}_2=\frac{-i \omega L}{r(R+r)-i \omega L(R+2 r)} \tilde{U}_1$$
\end{solcard}

\begin{solcard}{C4}{Hereinafter, work in the limit $r\ll R, \omega L$. Obtain expressions for the power of Joule losses $P_{r1}$ and $P_{r2}$ released in the primary and secondary windings. Express your answer in terms of the amplitude of the input voltage $U_0$, as well as $r$, $\omega$ and $L$. Calculate the total power of Joule losses on windings $P_{r} =P_{r1} + P_{r2}$ for a transformer, the parameters of which were given earlier in the problem. Provide detailed calculations for your calculations in the solution sheets.}{1.10}

Obtain expressions for the power of Joule losses $P_{r1}$ and $P_{r2}$ released in the primary and secondary windings. Express your answer in terms of the amplitude of the input voltage $U_0$, as well as $r$, $\omega$ and $L$.


Calculate the total power of Joule losses on windings $P_{r} =P_{r1} + P_{r2}$ for a transformer whose parameters are given earlier in the problem. Provide detailed calculations for your calculations in the solution sheets.


Let's write down the power released on the first and second windings: $$ P_{r1}=\frac{|\tilde{I_1}|^2r}{2}=\frac{U_0^2r}{2}\cdot \left| \frac{R+r-i \omega L}{r(R+r)-i \omega L(R+2 r)}\right|^2 $$$$ P_{r2}=\frac{|\tilde{I_2}|^2r}{2}=\frac{U_0^2r}{2}\cdot \left| \frac{-i \omega L}{r(R+r)-i \omega L(R+2 r)}\right|^2 $$Считая, что $r \ll R, \omega L$ we get:


Then the total power losses: $$ P_r = \frac{U_0^2r}{2}\cdot\frac{R^2+2\omega^2 L^2}{\omega^2L^2R^2} $$


Let's do the calculations: $$ U_0 = U_{\text{rms}}\cdot \sqrt{2}=311~\text{V}\\ \omega = 2\pi f = 314~ \frac{\text{rad}}{\mathrm{s}}\\ L=\frac{\mu\mu_0SN^2}{l} = 0.503~ \text{H}\\ $$length одного витка: $$ c= 2\pi \cdot \sqrt{\frac{S}{\pi}}=35.4 ~\text{cm}\\ $$Resistance обмотки: $$ r = \frac{1}{\sigma_\text{w}}\cdot\frac{Nc}{\frac{\pi d^2}{4}}=2.16~ \text{\Omega} $$then

\par\smallskip\noindent\textbf{Answer:} \textbackslash{}textbf{Answer:} $$ P_{r1}=\frac{U_0^2r}{2}\cdot\frac{R^2+\omega^2 L^2}{\omega^2L^2R^2} ,\qquad P_{r2}=\frac{U_0^2r}{2R^2} $$
\end{solcard}

\begin{solcard}{D1}{Write down the expression for the tangential electric field $E(r,t)$.}{0.30}

Let us write down the magnetic field flux through a circle of radius $r$:$$\Phi=B\pi r^2$$Закон электромагнитной индукции for круга radiusом $r$:$$\mathcal{E}=-\dot{\Phi}\\ 2\pi r E(r,t)=i\omega\pi r^2B_0 e^{-i\omega t} $$

\par\smallskip\noindent\textbf{Answer:} \textbackslash{}textbf{Answer:} $$ E(r, t) = \text{Re}\left(\frac{B_0\omega r}{2}e^{\frac{i\pi}{2}-i\omega t}\right)= \frac{B_0\omega r}{2} \sin{\omega t} $$
\end{solcard}

\begin{solcard}{D2}{Derive an expression for the average power of Joule losses $P_\mathrm{eddy}$ released in the core. Express your answer in terms of the conductivity of the core material $\sigma_\mathrm c$, its cross-sectional area $S$, length $\ell$, and $B_0$ and $f$.}{0.70}

Let us write down the power of Joule losses per unit volume: $$ p_\mathrm{eddy}(r, t) = \sigma_c E^2=\frac{\sigma_c B_0^2 \omega^2 r^2}{4}\sin^2\omega t $$averaging по времени we get: $$ \overline{p_\mathrm{eddy}} = \frac{\sigma_c B_0^2 \omega^2 r^2}{8} $$Then средняя power выделяющаяся на всем сердечнике: $$ P_\mathrm{eddy} = \frac{\sigma_c B_0^2 \omega^2 l}{8} \cdot \int_0^{\sqrt{\frac{S}{\pi}}}2\pi r^3 dr $$calculating

\par\smallskip\noindent\textbf{Answer:} \textbackslash{}textbf{Answer:} $$ P_\mathrm{eddy}  = \frac{\pi}{4}\cdot \sigma_c B_0^2 f^2 S^2 l $$
\end{solcard}

\begin{solcard}{D3}{Find the factor $A$.}{0.30}

Let us write the law of electromagnetic induction for a rectangle with dimensions $2z\times 2x$: $$ E_x(z, t)\cdot 4x = -\dot{\Phi} = 2x\cdot 2z\cdot \text{Re}~(i\omega B_0 e^{-i\omega t}) \\ E_x(z,t) =\text{Re}~(i\omega B_0 z e^{-i\omega t}) $$

\par\smallskip\noindent\textbf{Answer:} \textbackslash{}textbf{Answer:} $$ A = i\omega B_0 $$
\end{solcard}

\begin{solcard}{D4}{Write down an expression for the power $\mathrm dP/\mathrm dV$ released in such a plate per unit volume, averaged over the entire height of the plate, as well as over time. Responses may include $S$, $\ell$, $B_0$, $h$, $\sigma_{\mathrm c}$ and $f$. Let now the core be assembled from plates with the parameters indicated above in the problem. The load $R$ is still connected to the secondary winding. Express the power $P'_\mathrm{eddy}$ that would be released in the core due to Foucault currents, in terms of the parameters given in the problem. Calculate this power.}{1.40}

Let's write down the expression for instantaneous power at a certain point per unit area: $$\frac{\mathrm dP}{\mathrm dV}(z, t) = \sigma_c E^2(z, t) = \sigma_c \omega^2 B_0^2 z^2\sin^2{\omega t}$$averaging по времени we get: $$\frac{\mathrm dP}{\mathrm dV}(z) =\frac{1}{2} \sigma_c \omega^2 B_0^2 z^2$$теперь усредним по высоте $z$:$$ \frac{\mathrm dP}{\mathrm dV} =\frac{1}{2} \sigma_c \omega^2 B_0^2\cdot \frac{1}h\displaystyle\int\limits^{{h}/2}_{-{h}/2} z^2 dz = \frac{1}{24}\sigma_c \omega^2B_0^2h^2$$


Let's find an expression for the power allocated to the entire transformer: $$P'_{eddy}=\frac{\mathrm dP}{\mathrm dV}\cdot Sl = \frac{\pi^2}{6}\sigma_cB_0^2 f^2h^2Sl$$можно заметить, что here power пропорциональна $S$, not $S^2$, (как это было in the case of цельного сердечника), and therefore power потерь будет thereforeельно ниже. Теперь let us find амплитуду магнитного поля $B_0$. We use equation for связи поcurrentа магнитного поля via/through сердечник and напряжения на обмотке from части $A$: $$ U_1 = N\dot{\Phi} = NS(-i\omega B)\\ B_0= \frac{U_0}{2\pi fNS} = 0.99 ~\text{T} $$ рассчитаем $P'_\mathrm{eddy}$:


Now let's find the amplitude of the magnetic field $B_0$. We use the equation to relate the magnetic field flux through the core and the voltage on the winding from part $А$: $$ U_1 = N\dot{\Phi} = NS(-i\omega B)\\ B_0= \frac{U_0}{2\pi fNS} = 0.99 ~\text{T} $$ рассчитаем $P'_\mathrm{eddy}$:

\par\smallskip\noindent\textbf{Answer:} \textbackslash{}textbf{Answer:} $$ \frac{\mathrm dP}{\mathrm dV} = \frac{\pi^2}{6}\sigma_cB_0^2 f^2h^2 $$
\end{solcard}

\begin{solcard}{E1}{For simplicity, consider a transformer without load. Record the energy that the source imparts to the core when current $I$ flows through the primary winding and the flux through the core changes to $\mathrm d\Phi$. Then express this energy in terms of $H$, $\mathrm dB$ and core parameters.}{0.40}
\par\smallskip\noindent\textbf{Answer:} \textbackslash{}textbf{Answer:} $$ \mathrm dW =N\cdot I\,\mathrm d\Phi= N\frac{Hl}{N}S\,\mathrm dB=Sl\cdot H\,\mathrm dB $$
\end{solcard}

\begin{solcard}{E2}{We will assume that the last expression is valid in an arbitrary case. Write down the expression for the average power $P_{\mathrm{hyst}}$ hysteresis losses. Express your answer in terms of $H_\mathrm c$, $B_\mathrm s$ and the quantities given at the beginning of part C. Also get the numerical value.}{0.30}

For one period on the graph $B(H)$ one complete hysteresis loop will be described, which means $$\Delta W = Sl\cdot\Delta H \Delta B = Sl\cdot2H_c\cdot2B_s = 4SlH_cB_s$$Then средняя power энергии в следствии гистерезиса будет equals: $$P_\mathrm{hyst} = \frac{\Delta W}{T}$$ energy will be lost

\par\smallskip\noindent\textbf{Answer:} \textbackslash{}textbf{Answer:} $$ P_\mathrm{hyst} = 4SlH_cB_sf= 15~ \text{W} $$
\end{solcard}

\begin{solcard}{E3}{Assuming that all contributions calculated in parts C, D and E are small and independent, calculate the final efficiency of the transformer $\eta$.}{0.50}

$$\eta = 1-\frac{P_r+P_\mathrm{eddy}+P_\mathrm{hyst}}{P_0}$$

\par\smallskip\noindent\textbf{Answer:} \textbackslash{}textbf{Answer:} $$\eta = 91~\%$$
\end{solcard}

\end{document}